% Chapter 5 --- Substochastic kernels and the kernel-embedding anchor
% (paper §6, milestone M5).
%
% Mirrors theories/kernels/:
%   skern.v   the category Skern of substochastic kernels (mathcomp-
%             analysis fker-based, with the norm-≤-1 constraint)
%   kernel_embedding.v   paper Thm 6.5 — Skern → ICones is fully faithful
%
% Wave 2 is populated. The two headline-theorem statements
% (Thm 6.1 round-trips and Thm 6.5 fully faithfulness) are tagged
% \rocqok and \mathcompok: they are fully closed by Qed in the source
% tree and depend only on axioms from mathcomp.classical.boolp
% (propositional_extensionality, functional_extensionality_dep,
% constructive_indefinite_description) — no project-specific axioms.

%%%%%%%%%%%%%%%%%%%%%%%%%%%%%%%%%%%%%%%%%%%%%%%%%%%%%%%%%%%%%%%%%%%%%%%%
\chapter{Substochastic kernels and the kernel-embedding anchor}\label{ch:skern}

This is the headline chapter of the formalisation: the substochastic
kernel category $\Skern$ --- arguably the most natural source of
denotational semantics for first-order probabilistic programming ---
embeds fully and faithfully into $\ICones$ (paper Theorem~6.5). The
proof of the kernel-embedding anchor reduces, via paper Theorem~6.1 (the
natural isomorphism $\Path(X, B) \simeq \FMeas(X) \lin B$), to two
short corollaries of the round-trip identities established in
Chapter~\ref{ch:linhom}. All material in this chapter lives under
\texttt{theories/kernels/}, with the supporting machinery in
\texttt{theories/homs/bilin.v}.

\section{Paper Theorem 6.1: $\Path(X, B) \simeq \FMeas(X) \lin B$}%
\label{sec:thm61}

The starting point is the natural isomorphism between paths from an
$\Ar$-object $X$ to an integrable cone $B$ and integrable linear maps
from $\FMeas(X)$ to $B$. This is paper Theorem~6.1; we summarise the
construction and the two round-trip identities that drive
Theorem~6.5. The forward map is integration of the path against the
input measure; the inverse map is precomposition with the canonical
Dirac path.

\subsection{The Dirac path $\delta_X$}\label{sec:dirac-path}

\begin{definition}[Dirac measure as a finite measure]%
\label{def:dirac-fmeas}\rocqok\mathcompok
\rocq{Icones.homs.bilin.dirac_fmeas}%
\uses{def:fmeas-carrier}
For $X \in \Ar$ and $r \in X$, the Dirac mass $\delta_X(r)$ is the
finite measure on $X$ defined by
\[
  \delta_X(r)(U) =
    \begin{cases}
      1 & \text{if } r \in U,\\
      0 & \text{otherwise,}
    \end{cases}
  \qquad U \in \sigma_X.
\]
We package it as an element of $\FMeas(X)$ whose underlying measure
agrees with mathcomp-analysis's $\backslash\mathsf{d}_r$ on
measurable sets and is the canonical extension off the
$\sigma$-algebra. Its total mass is $\cnorm{\delta_X(r)} = 1$.
\end{definition}

\begin{lemma}[$\delta_X$ is a measurable path]\label{lem:dirac-path}%
\rocqok\mathcompok
\rocq{Icones.homs.bilin.dirac_path}%
\uses{def:dirac-fmeas, sec:path}
The assignment $r \mapsto \delta_X(r)$ is a bounded measurable path
$\delta_X \in \Path(X, \FMeas(X))$. Boundedness is immediate
($\cnorm{\delta_X(r)} = 1$); joint test-measurability against the
canonical test family of $\FMeas(X)$ --- the maps
$e_U \colon \mu \mapsto \mu(U)$ for $U \in \sigma_X$ --- reduces to
measurability of $r \mapsto \mathbf{1}_U(r)$, which is immediate.
\end{lemma}

\subsection{The forward map $I^B_X$}\label{sec:int-to-linhom}

\begin{definition}[Forward map of Theorem~6.1]%
\label{def:int-to-linhom}\rocqok\mathcompok
\rocq{Icones.homs.bilin.int_to_linhom}%
\uses{sec:icones, sec:linhom-icone, lem:dirac-path}
For $\beta \in \Path(X, B)$, define
\[
  I^B_X(\beta) : \FMeas(X) \lin B,
  \qquad I^B_X(\beta)(\mu) := \Pettis{X}{\beta\,d\mu}.
\]
The Rocq definition \rocq{Icones.homs.bilin.int_to_linhom_fun}
provides the underlying function; the full $\linhom$ packaging
exhibits its five defining fields:
\begin{enumerate}
  \item linearity in $\mu$
    (\rocq{Icones.homs.bilin.int_to_linhom_fun_linear}), paper
    Lemma~4.7 part bilinearity;
  \item $\omega$-continuity in $\mu$
    (\rocq{Icones.homs.bilin.int_to_linhom_fun_continuous}), paper
    Lemma~4.7 part $\omega$-continuity, via the
    rescaling-to-the-unit-ball argument and
    \texttt{integral\_omega\_cont\_meas};
  \item operator-norm boundedness with constant $\cnorm{\beta}_\Path$
    (\rocq{Icones.homs.bilin.int_to_linhom_fun_bounded}), from paper
    Lemma~4.2;
  \item measurable-path preservation in the $\FMeas$ argument
    (\rocq{Icones.homs.bilin.int_to_linhom_fun_pres_path}), paper
    Lemma~4.7 joint measurability;
  \item integral preservation
    (\rocq{Icones.homs.bilin.int_to_linhom_fun_pres_int}), via the
    cone-Fubini theorem and the scalar kernel-Tonelli identity
    \rocq{Icones.homs.bilin.icone_integral_kernel_tonelli} that
    repackages a measurable path of finite measures as a finite
    kernel and invokes mathcomp-analysis's \texttt{integral\_kcomp}.
\end{enumerate}
\end{definition}

\subsection{The inverse map $K^B_X$}\label{sec:linhom-to-int}

\begin{definition}[Inverse map of Theorem~6.1]%
\label{def:linhom-to-int}\rocqok\mathcompok
\rocq{Icones.homs.bilin.linhom_to_int}%
\uses{lem:dirac-path, sec:linhom-icone}
For $f : \FMeas(X) \lin B$, define
\[
  K^B_X(f) : \Path(X, B), \qquad
  K^B_X(f)(r) := f(\delta_X(r)).
\]
That is, $K^B_X(f) = f \circ \delta_X$. Boundedness follows from the
operator norm of $f$; measurability of the path follows from the
fact that $f$ preserves measurable paths and $\delta_X$ is one
(Lemma~\ref{lem:dirac-path}).
\end{definition}

\subsection{Round-trip $K \circ I = \id$}\label{sec:KI-roundtrip}

\begin{theorem}[Round-trip $K \circ I = \id$ at the path level]%
\label{thm:KI}\rocqok\mathcompok
\rocq{Icones.homs.bilin.K_I_int_to_linhom_path_E}%
\uses{def:int-to-linhom, def:linhom-to-int}
For every $\beta \in \Path(X, B)$,
\[
  K^B_X(I^B_X(\beta)) = \beta.
\]
\end{theorem}
\begin{proof}
By path extensionality
(\rocq{Icones.mcones.path.path_eq}). For $r \in X$,
\[
  K^B_X(I^B_X(\beta))(r)
   = I^B_X(\beta)(\delta_X(r))
   = \Pettis{X}{\beta\,d(\delta_X(r))}
   = \beta(r),
\]
where the last equality is the elementary identity
\rocq{Icones.homs.bilin.int_to_linhom_fun_dirac}: by
Pettis-uniqueness it suffices to check, for every test $m \in
\Mtest_B$, that
\[
  m(\beta(r))
    = \mathrm{fine}\!\int_X (m \circ \beta)(r')\,d(\delta_X(r))(r').
\]
The right-hand side agrees pointwise with the Lebesgue integral of
$m \circ \beta$ against $\backslash\mathsf{d}_r$ (since $\delta_X(r)$
agrees with $\backslash\mathsf{d}_r$ on measurable sets), and
mathcomp-analysis's \texttt{integral\_dirac} reduces this to
$m(\beta(r))$.
\end{proof}

\subsection{Round-trip $I \circ K = \id$}\label{sec:IK-roundtrip}

The reverse round-trip is more delicate because $\FMeas(X)$ is not
freely generated by Diracs at the cone level. The bridging lemma is
the \emph{Dirac approximation} of an arbitrary finite measure.

\begin{lemma}[Dirac approximation of finite measures]%
\label{lem:dirac-approx}\rocqok\mathcompok
\rocq{Icones.homs.bilin.icone_integral_dirac_path}%
\uses{lem:dirac-path, sec:pettis}
For every $\mu \in \FMeas(X)$,
\[
  \mu \;=\; \Pettis{X}{\delta_X\,d\mu} \;\in\; \FMeas(X).
\]
That is, every finite measure equals the Pettis integral of the
Dirac path against itself.
\end{lemma}
\begin{proof}
By measure extensionality (\rocq{Icones.mcones.fmeas.fmeas_eq}) we
check equality on every measurable set $U \in \sigma_X$. Applying the
defining Pettis equation of $\int \delta_X\,d\mu$ against the test
$e_U \in \Mtest_{\FMeas(X)}$ reduces the right-hand side to
\[
  \int_X \delta_X(r)(U)\,d\mu(r) \;=\; \int_X \mathbf{1}_U(r)\,d\mu(r) \;=\;
  \mu(U),
\]
where the last equality is \texttt{integral\_indic}.
\end{proof}

\begin{theorem}[Round-trip $I \circ K = \id$ at the linhom level]%
\label{thm:IK}\rocqok\mathcompok
\rocq{Icones.homs.bilin.I_K_int_to_linhom_E}%
\uses{def:int-to-linhom, def:linhom-to-int, lem:dirac-approx}
For every $f : \FMeas(X) \lin B$,
\[
  I^B_X(K^B_X(f)) = f.
\]
\end{theorem}
\begin{proof}
By $\linhom$-extensionality (\rocq{Icones.homs.linhom.linhom_eq}),
check pointwise on $\mu \in \FMeas(X)$. Unfolding,
\[
  I^B_X(K^B_X(f))(\mu)
    = \Pettis{X}{(r \mapsto f(\delta_X(r)))\,d\mu}
    = f\!\left(\Pettis{X}{\delta_X\,d\mu}\right)
    = f(\mu),
\]
where the middle step is the integral-preservation field of $f$
(applied to the measurable path $\delta_X$ from
Lemma~\ref{lem:dirac-path}), and the last is the Dirac
approximation, Lemma~\ref{lem:dirac-approx}.
\end{proof}

\begin{remark}[Why not package $I^B_X, K^B_X$ as an $\MCones$ iso?]
The two round-trips deliver a set-level bijection between
$\Path(X, B)$ and $\FMeas(X) \lin B$. Packaging $I^B_X$ as a full
$\mcones_{\mathrm{hom}}$ requires the pointwise norm-decrease
constraint \rocq{Icones.cones.cone_cat.cones_hom_norm_le1}
($\cnorm{f(x)} \le \cnorm{x}$), whereas $I^B_X(\beta)$ is only
bounded by $\cnorm{\beta}_\Path \cdot \cnorm{\mu}_\FMeas$ in general.
A normalisation wrapper or a relaxation of the bound is needed for
the full $\mcones$-iso packaging; we defer it. The kernel-embedding anchor
(Theorem~\ref{thm:thm65}) works directly with the
unit-norm-constrained version through $\Skern$, which is the
content of the next section.
\end{remark}

\section{The category $\Skern$}\label{sec:skern}

We now define the category of substochastic kernels following paper
\S6.1, with the unit-norm constraint baked into the morphism type.

\begin{definition}[Substochastic kernel]\label{def:skern-hom}%
\rocqok\mathcompok
\rocq{Icones.kernels.skern.Skern_hom}%
\uses{sec:path, def:fmeas-carrier}
For $X, Y \in \Ar$, a \emph{substochastic kernel} from $X$ to $Y$ is a
measurable path $\kappa \in \Path(X, \FMeas(Y))$ of operator norm
$\cnorm{\kappa}_\Path \le 1$. We write $\Skern(X, Y)$ for the set of
substochastic kernels, and pack each as a record
$\mathsf{MkSkernHom}(\kappa, h)$ with $h : \cnorm{\kappa}_\Path \le
1$. Equality of $\Skern$-morphisms reduces to equality of the
underlying paths
(\rocq{Icones.kernels.skern.Skern_hom_eq}) via
$\mathsf{Prop}$-irrelevance on the norm-bound proof.
\end{definition}

\begin{remark}
The paper's notation $\Skern(X, Y) = B_{\Path(X, \FMeas(Y))}$
(unit-ball paths) matches this definition exactly. Crucially we use
the project's own $\FMeas$ rather than mathcomp-analysis's
\texttt{\{finite\_measure\}} type, so that $\FMeas(Y)$ already
carries the $\ICone$ structure built in Chapter~\ref{ch:icones} and
all of Theorem~6.1 applies directly without further coercions.
\end{remark}

\begin{definition}[Identity in $\Skern$]\label{def:skern-id}%
\rocqok\mathcompok
\rocq{Icones.kernels.skern.Skern_id}%
\uses{lem:dirac-path, def:skern-hom}
For $X \in \Ar$, the identity at $X$ is the Dirac path:
$\id_X := \delta_X \in \Skern(X, X)$. The unit-norm bound
$\cnorm{\delta_X}_\Path \le 1$ holds with equality since each value
$\delta_X(r)$ has norm exactly $1$
(\rocq{Icones.kernels.skern.dirac_path_norm_le1}).
\end{definition}

\begin{definition}[Composition in $\Skern$]\label{def:skern-comp}%
\rocqok\mathcompok
\rocq{Icones.kernels.skern.Skern_comp}%
\uses{def:skern-hom, def:int-to-linhom}
Given $\lambda \in \Skern(X, Y)$ and $\kappa \in \Skern(Y, Z)$, the
Kleisli composite $\lambda \mathbin{;} \kappa \in \Skern(X, Z)$ is
defined by
\[
  (\lambda \mathbin{;} \kappa)(r)(U) := \int_Y \kappa(s)(U)\,d\lambda(r)(s)
  \quad (U \in \sigma_Z),
\]
or equivalently
\[
  (\lambda \mathbin{;} \kappa)(r) := I^{\FMeas(Z)}_Y(\kappa)(\lambda(r))
  \in \FMeas(Z).
\]
Path-preservation is provided by
\rocq{Icones.homs.bilin.int_to_linhom_fun_pres_path}, applied to
$\kappa$ and the measurability of $\lambda$. The norm bound
$\cnorm{\lambda \mathbin{;} \kappa}_\Path \le 1$ follows from the
$\Path$-level operator inequality
$\cnorm{\int \kappa\,d\nu}_\FMeas \le \cnorm{\kappa}_\Path
\cdot \cnorm{\nu}_\FMeas$ and the hypotheses
$\cnorm{\kappa}_\Path \le 1$, $\cnorm{\lambda(r)}_\FMeas \le 1$.
\end{definition}

\begin{proposition}[Category laws for $\Skern$]\label{prop:skern-cat}%
\rocqok\mathcompok
\uses{def:skern-id, def:skern-comp, thm:KI, lem:dirac-approx}
The data $(\Ar, \Skern_{\bullet,\bullet}, \id, \mathbin{;})$
satisfies the category axioms:
\begin{itemize}
  \item left identity
    \rocq{Icones.kernels.skern.Skern_compIl}: for $\lambda \in
    \Skern(X, Y)$, $\lambda \mathbin{;} \delta_Y = \lambda$;
  \item right identity
    \rocq{Icones.kernels.skern.Skern_compIr}: for $\lambda \in
    \Skern(X, Y)$, $\delta_X \mathbin{;} \lambda = \lambda$;
  \item associativity
    \rocq{Icones.kernels.skern.Skern_compA}: for composable
    $\kappa_1, \kappa_2, \kappa_3$,
    $\kappa_1 \mathbin{;} (\kappa_2 \mathbin{;} \kappa_3) =
     (\kappa_1 \mathbin{;} \kappa_2) \mathbin{;} \kappa_3$.
\end{itemize}
\end{proposition}
\begin{proof}
Both unit laws are pointwise statements: by path extensionality on
$r \in X$, we reduce to identities about $I^{\FMeas(Y)}_Y(\cdot)$
applied at a single point. The left identity becomes
$I^{\FMeas(Y)}_Y(\delta_Y)(\lambda(r)) = \lambda(r)$, which is
exactly the Dirac approximation (Lemma~\ref{lem:dirac-approx})
applied to $\mu = \lambda(r)$. The right identity becomes
$I^{\FMeas(Y)}_X(\lambda)(\delta_X(r)) = \lambda(r)$, which is
Theorem~\ref{thm:KI}'s elementary instance
\rocq{Icones.homs.bilin.int_to_linhom_fun_dirac}.
Associativity reduces to
$I^{\FMeas(X_4)}_{X_3}(\kappa_3) \circ I^{\FMeas(X_3)}_{X_2}(\kappa_2)
= I^{\FMeas(X_4)}_{X_2}(s \mapsto I^{\FMeas(X_4)}_{X_3}(\kappa_3)
(\kappa_2(s)))$ applied at $\kappa_1(r)$, which is the
integral-preservation field of $I^{\FMeas(X_4)}_{X_3}(\kappa_3)$,
i.e.\ \rocq{Icones.homs.bilin.int_to_linhom_fun_pres_int}.
\end{proof}

\section{The functor $\Skern \to \ICones$}\label{sec:skern-to-icones}

We now define the functor $\mathsf{Klin} : \Skern \to \ICones$ of paper
\S6.1, line~3058: it sends $X$ to $\FMeas(X)$, and a kernel $\kappa$
to its associated integrable linear map $I^{\FMeas(Y)}_X(\kappa)$.

\begin{definition}[$\mathsf{Klin}$ on objects]\label{def:Klin-obj}%
\rocqok\mathcompok
\rocq{Icones.kernels.kernel_embedding.Skern_to_ICones_obj}%
\uses{def:fmeas-carrier}
$\mathsf{Klin}(X) := \FMeas(X) \in \ICones$.
\end{definition}

\begin{definition}[$\mathsf{Klin}$ on morphisms]\label{def:Klin-mor}%
\rocqok\mathcompok
\rocq{Icones.kernels.kernel_embedding.Skern_to_ICones_mor}%
\uses{def:int-to-linhom, def:skern-hom}
For $\kappa \in \Skern(X, Y)$,
$\mathsf{Klin}(\kappa) := I^{\FMeas(Y)}_X(\kappa) \in
\ICones(\FMeas(X), \FMeas(Y))$. The full $\icones_{\mathrm{hom}}$
packaging re-exhibits the five fields of $I^{\FMeas(Y)}_X(\kappa)$
(linearity, $\omega$-continuity, boundedness, measurable-path
preservation, integral preservation) inherited from
Definition~\ref{def:int-to-linhom}, plus the pointwise
\rocq{Icones.cones.cone_cat.cones_hom_norm_le1} bound
\[
  \cnorm{\mathsf{Klin}(\kappa)(\mu)}_\FMeas \le \cnorm{\mu}_\FMeas,
\]
which uses $\cnorm{\kappa}_\Path \le 1$ together with
\rocq{Icones.icones.icone_integral.path_integral_norm_le}
to chain $\cnorm{\int \kappa\,d\mu}_\FMeas
\le \cnorm{\kappa}_\Path \cdot \cnorm{\mu}_\FMeas
\le 1 \cdot \cnorm{\mu}_\FMeas
= \cnorm{\mu}_\FMeas$ (see
\rocq{Icones.kernels.kernel_embedding.Skern_to_ICones_mor_norm_le1}).
The unfolding lemma
\rocq{Icones.kernels.kernel_embedding.Skern_to_ICones_mor_E} confirms
$\mathsf{Klin}(\kappa)(\mu) = \int_X \kappa\,d\mu$.
\end{definition}

\begin{proposition}[Functoriality of $\mathsf{Klin}$]\label{prop:Klin-fun}%
\rocqok\mathcompok
\uses{def:Klin-obj, def:Klin-mor, prop:skern-cat}
$\mathsf{Klin}$ preserves identity and composition:
\begin{itemize}
  \item \rocq{Icones.kernels.kernel_embedding.Skern_to_ICones_mor_id}:
    $\mathsf{Klin}(\delta_X) = \id_{\FMeas(X)}$ in $\ICones$;
  \item \rocq{Icones.kernels.kernel_embedding.Skern_to_ICones_mor_comp}:
    $\mathsf{Klin}(\kappa_1 \mathbin{;} \kappa_2)
    = \mathsf{Klin}(\kappa_2) \circ \mathsf{Klin}(\kappa_1)$.
\end{itemize}
\end{proposition}
\begin{proof}
The identity equation reduces, by $\icones_{\mathrm{hom}}$-extensional\-ity, to
$I^{\FMeas(X)}_X(\delta_X)(\mu) = \mu$, which is exactly the Dirac
approximation (Lemma~\ref{lem:dirac-approx}). The composition
equation reduces, again by extensionality on $\mu$, to the
integral-preservation identity for $I^{\FMeas(Z)}_Y(\kappa_2)$
applied to the measurable path $\kappa_1$
(\rocq{Icones.homs.bilin.int_to_linhom_fun_pres_int}).
\end{proof}

\section{Faithfulness}\label{sec:skern-faithful}

\begin{theorem}[Paper Theorem~6.5, faithful direction]%
\label{thm:thm65-faithful}\rocqok\mathcompok
\rocq{Icones.kernels.kernel_embedding.Skern_to_ICones_faithful}%
\uses{def:Klin-mor, thm:KI}
For $X, Y \in \Ar$ and $\kappa_1, \kappa_2 \in \Skern(X, Y)$, if
$\mathsf{Klin}(\kappa_1) = \mathsf{Klin}(\kappa_2)$ then $\kappa_1 = \kappa_2$.
\end{theorem}
\begin{proof}
Equality of $\icones_{\mathrm{hom}}$s implies pointwise equality of
the underlying functions:
$I^{\FMeas(Y)}_X(\kappa_1)(\mu) = I^{\FMeas(Y)}_X(\kappa_2)(\mu)$ for
every $\mu \in \FMeas(X)$. By
\rocq{Icones.homs.linhom.linhom_eq} this lifts to equality of the
$\linhom$-packages $I^{\FMeas(Y)}_X(\kappa_1)
= I^{\FMeas(Y)}_X(\kappa_2)$. Applying $K^{\FMeas(Y)}_X$ to both
sides and invoking the round-trip $K \circ I = \id$
(Theorem~\ref{thm:KI},
\rocq{Icones.homs.bilin.K_I_int_to_linhom_path_E}):
\[
  \kappa_1
   = K^{\FMeas(Y)}_X(I^{\FMeas(Y)}_X(\kappa_1))
   = K^{\FMeas(Y)}_X(I^{\FMeas(Y)}_X(\kappa_2))
   = \kappa_2,
\]
where the outer equalities hold at the underlying path level
$\Path(X, \FMeas(Y))$, and the lift to $\Skern_{\mathrm{hom}}$
follows by \rocq{Icones.kernels.skern.Skern_hom_eq}
($\mathsf{Prop}$-irrelevance on the norm-bound proof).
\end{proof}

\section{Fullness}\label{sec:skern-full}

\begin{theorem}[Paper Theorem~6.5, full direction]%
\label{thm:thm65-full}\rocqok\mathcompok
\rocq{Icones.kernels.kernel_embedding.Skern_to_ICones_full}%
\uses{def:Klin-mor, thm:IK, def:linhom-to-int}
For every $f \in \ICones(\FMeas(X), \FMeas(Y))$, there exists
$\kappa \in \Skern(X, Y)$ such that $\mathsf{Klin}(\kappa) = f$.
\end{theorem}
\begin{proof}
Set $\kappa := K^{\FMeas(Y)}_X(f) \in \Path(X, \FMeas(Y))$, i.e.\
$\kappa(r) := f(\delta_X(r))$. The unit-norm bound
$\cnorm{\kappa}_\Path \le 1$ follows pointwise: for every $r \in X$,
\[
  \cnorm{\kappa(r)}_\FMeas
   = \cnorm{f(\delta_X(r))}_\FMeas
   \le \cnorm{\delta_X(r)}_\FMeas
   = 1,
\]
where the inequality is \rocq{Icones.cones.cone_cat.cones_hom_norm_le1}
for $f$ (in its $\icones_{\mathrm{hom}}$ packaging) and the last
equality is \rocq{Icones.homs.bilin.dirac_fmeas_norm} (see
\rocq{Icones.kernels.kernel_embedding.icones_to_skern_norm_le1}). So $\kappa$
is a substochastic kernel.

Now $f$ itself, viewed as a $\linhom$ via the helper
\rocq{Icones.kernels.kernel_embedding.icones_to_linhom} (which repackages the
$\icones_{\mathrm{hom}}$ fields with operator-norm constant
$M := 1$), satisfies the round-trip $I \circ K = \id$
(Theorem~\ref{thm:IK},
\rocq{Icones.homs.bilin.I_K_int_to_linhom_E}):
\[
  I^{\FMeas(Y)}_X(\kappa)
   = I^{\FMeas(Y)}_X(K^{\FMeas(Y)}_X(f))
   = f.
\]
Extracting underlying functions on both sides
(by $\icones_{\mathrm{hom}}$-extensionality
\rocq{Icones.icones.icone_cat.icones_hom_eq}) gives
$\mathsf{Klin}(\kappa) = f$ in $\ICones$, see
\rocq{Icones.kernels.kernel_embedding.Skern_to_ICones_mor_to_skern}.
\end{proof}

\section{Theorem 6.5 --- the kernel-embedding anchor}\label{sec:thm65}

\begin{theorem}[Paper Theorem~6.5, the kernel-embedding anchor]%
\label{thm:thm65}\rocqok\mathcompok
\rocq{Icones.kernels.kernel_embedding.Skern_to_ICones_fully_faithful}%
\uses{thm:thm65-faithful, thm:thm65-full}
For all $X, Y \in \Ar$, the assignment
\[
  \mathsf{Klin}_{X,Y} : \Skern(X, Y) \to \ICones(\FMeas(X), \FMeas(Y)),
  \qquad \kappa \mapsto I^{\FMeas(Y)}_X(\kappa)
\]
is a bijection. Equivalently, the functor
$\mathsf{Klin} : \Skern \to \ICones$ defined by
$\mathsf{Klin}(X) := \FMeas(X)$,
$\mathsf{Klin}(\kappa) := I^{\FMeas(Y)}_X(\kappa)$ is full and faithful.
\end{theorem}
\begin{proof}
Combine Theorems~\ref{thm:thm65-faithful} and~\ref{thm:thm65-full}.
\end{proof}

\section{Design choices and the axiom budget}\label{sec:skern-design}

We close with a few remarks on the formalisation choices behind
Theorem~\ref{thm:thm65}, and on the axiom budget of the development.

\paragraph{$\FMeas$ vs.\ \texttt{\{finite\_measure\}}.}
Throughout we use the project's $\FMeas(X)$ record
(\rocq{Icones.mcones.fmeas.fmeas}), not mathcomp-analysis's
\texttt{\{finite\_measure\}} type, because $\FMeas(X)$ already
carries the cone, measurable-cone, and integrable-cone structure
built in Chapters~\ref{ch:mcones} and~\ref{ch:icones}. The
canonicality invariant of $\FMeas$ --- the underlying $\bar
R$-valued additive function is canonically extended off the
$\sigma$-algebra --- gives definitional equality on the carrier and
keeps $\FMeas(X)$ a member of $\ICones$ on the nose.

\paragraph{The Dirac approximation as the bridging lemma.}
Lemma~\ref{lem:dirac-approx},
\rocq{Icones.homs.bilin.icone_integral_dirac_path}, is the
mathematical content that makes Theorem~6.1's $I \circ K = \id$ go
through. It is the cone-level translation of the standard measure-
theoretic identity $\mu = \int \delta_r\,d\mu(r)$. The Rocq proof
checks the equation on every measurable $U$ via the Pettis equation
against the test $e_U \in \Mtest_{\FMeas(X)}$ and reduces to
mathcomp-analysis's \texttt{integral\_dirac} at the
$\sigma$-algebra level.

\paragraph{Why faithfulness and fullness are split.}
Each direction of Theorem~\ref{thm:thm65} is one short corollary of
one round-trip identity (faithful from $K \circ I = \id$, full from
$I \circ K = \id$). Splitting them gives reusable lemmas
\texttt{Skern\_to\_ICones\_faithful} and
\texttt{Skern\_to\_ICones\_full} that downstream developments
can call directly, while \texttt{Skern\_to\_ICones\_fully\_faithful}
just bundles them for the headline statement.

\paragraph{Why $I^B_X$ is not packaged as an $\MCones$-iso.}
Packaging the natural iso $\Path(X, B) \simeq \FMeas(X) \lin B$ as an
$\mcones_{\mathrm{iso}}$ requires either a normalisation wrapper on
$I^B_X(\beta)$ (whose operator norm is $\cnorm{\beta}_\Path$, not
$1$) or a relaxation of the $\mcones_{\mathrm{hom}}$ pointwise
norm-decrease constraint. We sidestep the difficulty for the anchor by
stating Theorem~\ref{thm:thm65} directly for $\Skern_{\mathrm{hom}}$,
which already carries the unit-norm constraint. The full
$\mcones$-iso packaging and naturality in $X$ and $B$ are not needed
for the kernel-embedding anchor.

\paragraph{Axiom budget.}
The kernel-embedding anchor,
\rocq{Icones.kernels.kernel_embedding.Skern_to_ICones_fully_faithful}, is
axiom-clean in the following sense: a \texttt{Print Assumptions}
call on the theorem reports only the three classical axioms exported
by \texttt{mathcomp.classical.boolp} and used pervasively by
\mathcompA{}:
\begin{itemize}
  \item \texttt{propositional\_extensionality},
  \item \texttt{functional\_extensionality\_dep},
  \item \texttt{constructive\_indefinite\_description}.
\end{itemize}
No project-specific axiom is introduced. This is why both
Theorem~\ref{thm:thm65} and all of its lemmas are tagged
\rocqok\ and \mathcompok\ in this chapter.

This is also the development's \emph{regression anchor}: every
downstream wave (the tensor $\otimes$ and SMCC of
Chapter~\ref{ch:tensor}, the stable category $\mathbf{SCones}$ of
Chapter~\ref{ch:stable}, the exponential $\mathord{!}$ and the Seely
category of Chapter~\ref{ch:exponential}, and the beyond-the-paper CBV
model of Chapter~\ref{ch:cbv}) leaves the
\texttt{Skern\_to\_ICones\_fully\_faithful} axiom budget unchanged ---
verified by a \texttt{Print Assumptions} after each merge.

\paragraph{The kernel-embedding anchor and the rest of the development.}
Theorem~\ref{thm:thm65} delivers the M5 milestone of
\texttt{PLAN.md}: paper \S2 (M1, $\Cones$), \S3 (M2, $\MCones$,
$\FMeas$, $\Path$), \S4 (M3, $\ICones$, Pettis, cone-Fubini), \S5.1
+ \S5.2 (M4, the internal hom $\lin$), and \S6 (M5, $\Skern$,
Theorem~6.5) are all formalised. The full linear-logic model of
paper \S5.3--\S5.5, \S7 and \S9 has since landed axiom-free
(Chapters~\ref{ch:tensor}, \ref{ch:stable}, \ref{ch:exponential}); the
beyond-the-paper CBV model of Chapter~\ref{ch:cbv} is also axiom-free.
What remains genuinely open is paper \S8 (analytic functions, the
category $\mathrm{ACONES}$), the \S9 Eilenberg--Moore full-subcategory
theorem ($\mathrm{ARCAT}$ on Polish / standard-Borel spaces embeds into
$\mathrm{EM}(\mathord{!})$ as a full subcategory --- needing a
Polish / standard-Borel layer not yet formalised), and the \S10 PCS
embedding.
